\begingroup
\begin{tikzpicture}[
    >=Stealth,
    font=\sffamily,
    token/.style={rounded corners=2pt, draw=black!25, fill=plotBlue!8,
        minimum width=1.75cm, minimum height=0.62cm, align=center,
        font=\small},
    emb/.style={rounded corners=2pt, draw=plotBlue!45, fill=plotBlue!10,
        minimum width=1.55cm, minimum height=0.50cm},
    transformer/.style={rounded corners=6pt, draw=plotAccent!60,
        fill=plotAccent!5, line width=0.9pt, minimum width=10.6cm,
        minimum height=4.2cm},
    subblock/.style={rounded corners=4pt, draw=plotAccent!55, fill=white,
        line width=0.75pt, minimum width=3.2cm, minimum height=1.05cm,
        align=center, font=\small},
    stage/.style={rounded corners=5pt, draw=plotBlue!45, fill=plotBlue!5,
        line width=0.75pt, minimum width=4.0cm, minimum height=1.0cm,
        align=center},
    arrow/.style={->, draw=black!60, line width=0.75pt},
    flow/.style={->, draw=plotBlue!75, line width=0.9pt},
    small/.style={font=\scriptsize, text=black!70, align=center},
    heading/.style={font=\small\bfseries, align=center}
]

\path[use as bounding box] (-6.0,-6.35) rectangle (6.0,11.15);

% Tokeny wejściowe.
\node[heading] at (0,10.70) {Tokeny};
\foreach \word/\name/\x in {
    Lorem/tok1/-4.0,
    ipsum/tok2/-2.0,
    dolor/tok3/0,
    sit/tok4/2.0,
    amet/tok5/4.0
} {
    \node[token] (\name) at (\x,9.95) {\texttt{\word}};
}

% Reprezentacje tokenów wraz z informacją o pozycji.
\foreach \name/\x in {emb1/-4.0,emb2/-2.0,emb3/0,emb4/2.0,emb5/4.0} {
    \node[emb] (\name) at (\x,8.05) {};
    \foreach \dx in {-0.48,-0.24,0,0.24,0.48} {
        \draw[plotBlue!45, line width=0.35pt]
            ($(\name.center)+(\dx,-0.21)$) -- ++(0,0.42);
    }
}
\foreach \a/\b in {tok1/emb1,tok2/emb2,tok3/emb3,tok4/emb4,tok5/emb5} {
    \draw[arrow] (\a.south) -- (\b.north);
}
\node[heading, fill=white, inner xsep=7pt, inner ysep=2pt] at (0,8.95)
    {Wektory tokenów + pozycja};

% Stos bloków Transformera.
\node[transformer, fill=plotBlue!6] (tr4) at (0.48,4.47) {};
\node[transformer, fill=plotBlue!9] (tr3) at (0.32,4.63) {};
\node[transformer, fill=plotBlue!12] (tr2) at (0.16,4.79) {};
\node[transformer] (tr1) at (0,4.95) {};

\draw[flow] (emb3.south) -- (tr1.north);
\node[font=\bfseries, text=plotAccent!85] at (0,6.41)
    {Stos bloków Transformera};
\node[small] at (0,6.02)
    {każdy blok zawiera mechanizm attention oraz sieć feed-forward};

\node[subblock, fill=plotAccent!11] (attn3) at (-2.04,4.47) {};
\node[subblock, fill=plotAccent!16] (attn2) at (-2.20,4.63) {};
\node[subblock] (attn1) at (-2.36,4.79)
    {multi-head\\attention};
\node[subblock] (ffn) at (2.28,4.79)
    {feed-forward\\network};
\draw[flow] (attn1.east) -- (ffn.west);

% Warstwa wyjściowa.
\node[stage] (output) at (0,1.45) {warstwa liniowa + softmax};
\draw[flow] (tr1.south) -- (output.north);

% Rozkład prawdopodobieństwa następnego tokenu.
\node[heading, fill=white, inner xsep=6pt, inner ysep=2pt] at (0,0.15)
    {Prawdopodobieństwa następnego tokenu};
\draw[arrow] (output.south) -- (0,0.48);

\node[rounded corners=5pt, draw=black!20, fill=black!1,
    minimum width=10.6cm, minimum height=4.35cm] (probbox) at (0,-2.45) {};

\foreach \tok/\pct/\w/\y in {
    consectetur/42\%/3.35/-1.05,
    elit/27\%/2.15/-1.75,
    adipiscing/16\%/1.28/-2.45,
    sed/8\%/0.64/-3.15,
    do/7\%/0.56/-3.85
} {
    \node[font=\scriptsize, anchor=east] at (-2.25,\y) {\texttt{\tok}};
    \draw[plotGrid, line width=3.0pt] (-1.95,\y) -- (2.70,\y);
    \draw[plotAccent, line width=3.0pt] (-1.95,\y) -- ++(\w,0);
    \node[font=\scriptsize, anchor=west, text=black!70] at (2.95,\y) {\pct};
}

\node[small] at (0,-5.15)
    {Warstwa wyjściowa wyznacza rozkład nad słownikiem modelu.};

\end{tikzpicture}
\endgroup
